\section{The Sobolev--Hilbert Trace}

The residue cannot enter completion as a naked symbol. Completion is earned
only when the grammar has a norm that reads activity, detects zero, and
controls Cauchy refinement. Earlier chapters grounded that norm in the
positive invariant: energy squared became a readable activity, and the
square-root reading opened the Hilbert door. Here the same door must receive
the extensional Cohen ledger and the successor-closed commutator residue.

The finite norm has a simple obligation. It must read zero exactly when the
residue has no activity left to carry. In schematic form,
\[
  \|r\| = 0 \quad\hbox{only when}\quad r \hbox{ is the zero residue.}
\]
The formula is not a new foundation. It records the old boundary-anchored
positivity in the language needed for completion. A norm that allowed
nonzero residue to pass as zero would reopen the null drift that the
concrete operator had already killed.

The pre-Hilbert skeleton is built from the countable extensional ledger.
Its elements are finite-supported readings of the obligations already
carried: atoms, bounds, decisions, and residues. Addition and scaling are
allowed only insofar as they preserve the finite compatibility discipline.
The skeleton is not yet the completed space. It is the countable stage on
which Cauchy refinement can be read.

Completion begins when sequences in this skeleton become Cauchy under the
norm. A sequence is Cauchy when late refinements no longer differ by more
than the chosen bound. This is the same finite approach that has guided the
book from the beginning: the grammar does not claim to have traversed every
point, but it can say that every later admissible refinement is being squeezed
within the bound already named. The Hilbert completion is the quotient of
such Cauchy refinement histories by unreadable late-stage difference.

The finite residue maps into this completion by remembering its origin. A
completed residue is not a residue without ancestry. It is the completed
reading of a finite residual history whose stages remain recoverable from
the ledger. The map into completion therefore carries the residue; it does
not dissolve it. If a completed object cannot be traced back to finite
residual obligations, it is not admitted by this grammar.

This is the first sense of trace in the section. A trace is the readable
memory of the finite source inside the completed reading. The trace does not
make the completion smaller. It makes it answerable. When a completed
residue is challenged, the grammar can ask which finite supports, bounds,
and successor comparisons forced it. If none can be given, the alleged
completion is optional structure, not a forced reading.

The second sense of trace prepares the boundary. Sobolev language is useful
because it knows how an interior activity can have a boundary reading. But
the chapter must use that language only after the finite discipline is in
place. The boundary trace is not a gift from smooth analysis. It will be the
completed reading of terminal finite residue. The present section builds the
door through which that terminal residue can pass.

Most important, the geometry/gauge split survives the passage. The geometric
certificate family carries boundary separation, mesh refinement, and
curvature-like residue. The gauge certificate family carries transported
orientation, charge, and commutator residue. Completion may complete both
families, but it may not identify their channels. The finite split lifts as
orthogonality in the Sobolev--Hilbert reading.

In symbols, the discipline has the shape
\[
  G \perp Q,
\]
where \(G\) names the completed geometric channel and \(Q\) names the
completed gauge channel. The symbol \(\perp\) is again only a record. It
says that the completed inner product respects the finite carrier split. It
does not say that geometry and gauge are unrelated. It says that their
relation is not an interior scalar product between two freely summable
components.

Precision again illustrates the distinction. A spline completion may supply
intermediate values that were not observed, but those values remain
consequences of the anchor grammar. They do not become independent
measurements. Likewise, the Sobolev--Hilbert completion supplies completed
residue readings, but those readings remain consequences of finite
conditions, successor residue, and the norm that bounds refinement.

The section therefore lifts finite residue without betraying finitude. The
countable extensional ledger enters a pre-Hilbert skeleton. The norm detects
zero activity and reads Cauchy refinement. The finite residue maps into
completion with its trace intact. The geometric and gauge certificate
families complete as split channels. Once that split has lifted, the next
attempted shortcut is visible: an interior mixed scalar in the continuum.
The grammar must remove it.
